\chapter{Constrained Optimization} \label{ch:constrained-optimization}

Constrained optimization is the problem of finding a minimum (or maximum) of a function $f(x)$ called the
\say{objective function}, subject to a number of constraints of the following types:
\bit
\item $h_{i}(x) =0, \:\:\: i=1,2,\ldots$ called \say{equality constraints}
\item $g_{i}(x) \leq 0, \:\:\: i=1,2,\ldots$ called \say{inequality constraints}
\eit

Let's start first with the equality constraints and the we will add inequality constraints.

\section{Equality Constrained Optimization}

The formulation of the optimization problem is to optimize $f(x)$ subject to $h_{i}(x) = 0, \:\:\: i=1,2,\ldots$.
Let's assume for simplicity only one constraint $h_{1}(x) = h(x) = 0$. The idea here is that the point that $f(x)$
touches $h(x)$ is the point that $f(x)$ is minimum (or maximum) while the constraint is also valid. At that point $f
(x)$ is parallel to $h(x)$ and the tangents $\nabla_{\boldsymbol{w}} f(x)$ and $\nabla_{\boldsymbol{w}} h(x)$ which
are perpendicular to $f(x)$ and $h(x)$ respectively, are also parallel to each other. Hence, since
$\nabla_{\boldsymbol{w}} f(x)$ and $\nabla_{\boldsymbol{w}} h(x)$ are parallel this translates to:
\bse
\nabla_{\boldsymbol{w}} f(x) = \mu \nabla_{\boldsymbol{w}} h(x)
\ese

\vspace{3pt}

which is the condition for the $f(x)$ to be minimum (or maximum) while $h(x) = 0$. \v

Without loss of generality the condition for many constraints reads:
\bse
\nabla_{\boldsymbol{w}} f(x) = \sum_{i} \mu_{i} \nabla_{\boldsymbol{w}} h_{i}(x)
\ese

or by bringing everything in one side:
\begin{align*}
& \nabla_{\boldsymbol{w}} f(x) - \sum_{i} \mu_{i} \nabla_{\boldsymbol{w}} h_{i}(x) = 0 \\
& \nabla_{\boldsymbol{w}} (f(x) - \sum_{i} \mu_{i} h_{i}(x)) = 0
\end{align*}

At this point we define the \say{Lagrangian} as follows:
\bse
\mathcal{L}(x, \mu_{i}) = f(x) + \sum_{i} \mu_{i} h_{i}(x)
\ese

where $\mu_{i}$ are called \say{Lagrange multipliers}. Subsequently the necessary conditions for optimization of $
\mathcal{L}$ turns to:
\bse
\nabla_{\boldsymbol{w}} \mathcal{L} = 0 \:\:\: \text{and} \:\:\: \frac{\partial \mathcal{L}}{\partial \mu_{i}} = 0
\ese

The solution of this system of equations minimizes (or maximizes) $f(x)$ subject to $h_{i}(x) = 0, \:\:\: \forall i$.

\section{Equality \& Inequality Constrained Optimization}

Now on top of equality constraints we also have inequality constraints. The formulation of the optimization problem
is to optimize $f(x)$ subject to $h_{i}(x) = 0, \:\:\: i=1,2,\ldots$ and $g_{i}(x) \leq 0, \:\:\: i=1,2,\ldots$.
Following a similar way of thinking as before, although a bit more technical, we can show (but we won't) that if the
following four conditions, called \say{Karush - Kuhn - Talker conditions} (KKT), hold:
\bit
\item $h_{i}(x) = 0, \:\:\: \forall i$
\item $g_{i}(x) \leq 0, \:\:\: \forall i$
\item $\lambda_{i} \leq 0, \:\:\: \forall i$
\item $\lambda_{i} g_{i}(x) = 0, \:\:\: \forall i$
\eit

then there exist constants $\mu_{i}$ and $\lambda_{i}$ called \say{KKT multipliers} such that:
\bse
\nabla_{\boldsymbol{w}} f(x) =
\sum_{i} \mu_{i} \nabla_{\boldsymbol{w}} h_{i} (x) + \sum_{i} \lambda_{i} \nabla_{\boldsymbol{w}} g_{i}(x)
\ese

By following the same philosophy as for the equality constrained optimization we define the \say{Lagrangian} as follows:
\bse
\mathcal{L}(x, \mu_{i}, \lambda_{i}) = f(x) + \sum_{i} \mu_{i} h_{i} (x) + \sum_{i} \lambda_{i} g_{i}(x)
\ese

and the necessary conditions for optimization of $ \mathcal{L}$ turn to:
\bse
\nabla_{\boldsymbol{w}} \mathcal{L}
= 0 \:\:\: \text{and} \:\:\: \frac{\partial \mathcal{L}}{\partial \mu_{i}}
= 0 \:\:\: \text{and} \:\:\: \frac{\partial \mathcal{L}}{\partial \lambda_{i}}
= 0
\ese

This final form of the optimization problem is usually called the \say{dual optimization problem}. \v

This is the most general case of constrained optimization. If there are no equality constraints $h_{i}(x)$ then we
simply have a theory for inequality constrained optimization. If there are no inequality constraints $g_{i}(x)$ then
the whole theory turns to the equality constrained optimization problem we developed previously and the KKT
multipliers turn to Lagrangian multipliers. Finally, if there are no equality constraints $h_{i}(x)$ neither
inequality constraints $g_{i}(x)$ the theory is just a usual optimization problem where we just find the solution
where the derivative of $f(x)$ is zero.

\chapter{Kernels} \label{ch:kernels}

\bd [Kernel]
Let $\bar{x}$ and $\bar{x}^\prime$ be two vectors of space $X$ and $\Phi$ a non-linear transformation. We define
$\bar{z}$ and $\bar{z}^\prime$ as the transformed vectors $\Phi(\bar{x})$ and $\Phi(\bar{x}^\prime)$:
\begin{align*}
& \bar{x} \in X \xrightarrow{\text{$\Phi$}} \bar{z} = \Phi(\bar{x}) \in Z \\
& \bar{x}^\prime \in X \xrightarrow{\text{$\Phi$}} \bar{z}^\prime = \Phi (\bar{x}^\prime) \in Z
\end{align*}

We define the \textbf{kernel} of space $Z$ as the function that is equal to the inner product of the transformation
vectors:
\bse
K(\bar{x}, \bar{x}^\prime) = \bar{z}^T \bar{z}^\prime
\ese
\ed

Let's for example assume the following non-linear transformation:
\begin{align*}
& \bar{x} = (x_1, x_2) \xrightarrow{\text{$\Phi$}} \bar{z} = \Phi (\bar{x}) = (1, x_1^2, x_2^2, \sqrt{2} x_1,
\sqrt{2} x_2, \sqrt{2} x_1 x_2) \\
& \bar{x}^\prime = (x_1^\prime, x_2^\prime) \xrightarrow{\text{$\Phi$}} \bar{z}^\prime = \Phi(\bar{x}^\prime) =
(1, x_1^{2^\prime}, x_2^{2^\prime}, \sqrt{2} x_1^\prime, \sqrt{2} x_2^\prime, \sqrt{2} x_1^\prime x_2^\prime)
\end{align*}

Then for the inner product:
\bse
\bar{z}^T \bar{z}^\prime = 1 + x_1^2 x_1^{2^\prime} + x_2^2 x_2^{2^\prime} + 2 x_1 x_1^\prime + 2 x_2 x_2^\prime + 2
x_1 x_2 x_1^\prime x_2^\prime
\ese

However we can get to the same result by simply defining a Kernel of the form:
\begin{align*}
K(\bar{x}, \bar{x}^\prime) &= (1 + \bar{x} \bar{x}^\prime)^2 \\
&= (1 + x_1 x_1^\prime + x_2 x_2^\prime )^2 \\
&= 1 + x_1^2 x_1^{2^\prime} + x_2^2 x_2^{2^\prime} + 2 x_1 x_1^\prime + 2 x_2 x_2^\prime
+ 2 x_1 x_2 x_1^\prime x_2^\prime
\end{align*}

Hence by knowing the Kernel of a space $Z$ of some non-linear transformation $\Phi$ we can compute inner products
without the need of transforming vectors from $X$ to $Z$. \v

The kernel trick is to use this idea in the opposite direction. Namely, to assume that a function $K(\bar{x},
\bar{x}^\prime)$ is the kernel of some space $Z$ for some non-linear transformation $\Phi$ and to compute inner
products without even knowing the transformation. \v

The question that arises is how do we know that some function $K(\bar{x}, \bar{x}^\prime)$ is actually the kernel of
a space $Z$. There are three approaches to this problem:
\ben
\item By construction (as we did in the example above).
\item By Mercer's condition that states that $K(\bar{x}, \bar{x}^\prime)$ is a valid kernel for some space $Z$ if
\bse
\int K(\bar{x}, \bar{x}^\prime) g(\bar{x}) g(\bar{x}^\prime) d\bar{x} d\bar{x}^\prime \geq 0 \:\:\: \forall \:\:\:
\text{square integrable functions} \:\:\: g(\bar{x})
\ese

\item Sometimes we don't care if $K(\bar{x}, \bar{x}^\prime)$ is a valid kernel for some space $Z$ as long as it does
the job.
\een

\chapter{Applied Machine Learning}

In this section we will cover the applied side of machine learning. This checklist can guide you through your machine
learning projects. There are eight main steps:
\ben
\item \textbf{Frame The Problem}
\bit
\item Define the objective in business terms.
\item How will your solution be used?
\item What are the current solutions/workarounds (if any)?
\item How should you frame this problem (supervised/unsupervised, online/offline, etc)?
\item How should performance be measured?
\item Is the performance measure aligned with the business objective?
\item What would be the minimum performance needed to reach the business objective?
\item What are comparable problems? Can you reuse experience or tools?
\item Is human expertise available? How would you solve the problem manually?
\item List the assumptions you (or others) have made so far and verify them if possible.
\eit

\item \textbf{Get The Data}
\bit
\item List the data you need and how much you need.
\item Find and document where you can get that data.
\item Check how much space it will take.
\item Check legal obligations, and get access authorization if necessary.
\item Create a workspace (with enough storage space).
\item Get the data.
\item Make sure you data are not of insufficient quantity, poor quality, or non representative of the population.
Especially validation and test set must be as representative as possible..
\item Convert the data to a format you can easily manipulate (without changing the data itself).
\item Ensure sensitive information is deleted or protected (e.g anonymized).
\item Make a copy of the original dataset and work with this one.
\item Automate as much as possible so you can easily get fresh data.
\eit

\item \textbf{Explore The Data}
\bit
\item Create a Jupyter notebook to keep a record of your data exploration.
\item Study each attribute and its characteristics: name, type, percentage of missing values, noisiness and type of
noise, outliers, usefulness for the task basic statistics (mean, standard deviation, type of distribution, etc).
\item For supervised learning tasks, identify the target attribute(s).
\item Sample a test set, put it aside, and never look at it (stratify if needed).
\item Visualize the data.
\item Study the correlations between attributes.
\item Study how you would solve the problem manually.
\item Identify extra data that would be useful.
\item Document what you have learned.
\item Try to get insights from a field expert for these steps.
\eit

\item \textbf{Prepare The Data}
\bit
\item Write functions for all data transformations you apply, for the following reasons:
\bit
\item So you can easily prepare the data the next time you get a fresh dataset.
\item So you can apply these transformations in future projects.
\item To clean and prepare the test set.
\item To clean and prepare new data instances once your solution is live.
\item To make it easy to treat your preparation choices as hyperparameters.
\eit

\item Clean the data: fix or remove outliers, fill in missing values (e.g with zero, mean, median,etc) or drop
their rows (or columns).
\item Feature selection/engineering:
\bit
\item Drop the attributes that provide no useful information for the task.
\item Drop variables that have a very high percentage of missing values.
\item Drop variables that have a very low variation (i.e.\ not too much information).
\item Drop variables that have very low correlation with the target.
\item Find variables that are highly correlated with each other (i.e.\ same behaviour), and keep the ones that have
higher correlation with the target (drop the other ones).
\item Identify the promising transformations you may want to apply (e.g combine existing features to produce more
useful ones or add promising transformations of features).
\item Select best features based on a metric. Either start with one variable and add more (forward selection), or
start with all variables and eliminate (backward elimination or recursive feature elimination).
\item Discretize continuous features.
\item Decompose features (e.g categorical, date/time, etc).
\item Standardize or normalize features (feature scaling). Machine learning algorithms don't perform well when the
input numerical attributes have very different scales. There are two common ways to get all attributes to have the
same scale: normalization and standardization. In normalization the values are shifted and rescaled so that they end
up ranging from 0 to 1. We do this by subtracting the minimum value and dividing by the maximum minus the minimum.
Standardization first subtracts the mean value (so standardized values always have a zero mean), and then divides by
the standard deviation so that the resulting distribution has unit variance. Unlike normalization, standardization
does not bound values to a specific range, which may be a problem for some algorithms. However, standardization is
much less affected by outliers.
\eit
\eit

\item \textbf{Shortlist Promising Models}
\bit
\item If the data is huge, you may want to sample smaller training sets so you can train many different models in a
reasonable time (be aware that this penalizes complex models such as large neural nets or Random Forests).
\item Train many quick-and-dirty models from different categories using standard parameters.
\item Measure and compare their performance.
\item For each model, use K-fold cross-validation and compute the mean and standard deviation of the performance
measure on the K folds.
\item Analyze the most significant variables for each algorithm.
\item Analyze the types of errors the models make. What data would a human have used to avoid these errors?
\item Shortlist the top three to five most promising models, preferring models that make different types of errors.
\eit

\item \textbf{Fine-Tune The System}
\bit
\item You will want to use as much data as possible for this step, especially as you move toward the end of fine-tuning.
\item Fine-tune the hyperparameters using cross-validation.
\item Try ensemble methods. Combining your best models will often produce better performance than running them
individually.
\item Once you are confident about your final model, measure its performance on the test set to estimate the
generalization error.
\item Don't tweak your model after measuring the generalization error since you would just start overfitting the test
set.
\eit

\item \textbf{Present Your Solution}
\bit
\item Document what you have done.
\item Create a nice presentation.
\item Make sure you highlight the big picture first.
\item Explain why your solution achieves the business objective.
\item Don't forget to present interesting points you noticed along the way.
\item Describe what worked and what did not.
\item List your assumptions and your system's limitations.
\item Ensure your key findings are communicated through beautiful visualizations or easy-to-remember statements.
\eit

\item \textbf{Launch}
\bit
\item Get your solution ready for production.
\item Write monitoring code to check your system's live performance at regular intervals and trigger alerts when it
drops.
\item Beware of slow degradation: models tend to ``rot'' as data evolves.
\item Measuring performance may require a human pipeline.
\item Also monitor your inputs' quality.
\item Retrain your models on a regular basis on fresh data (automate as much as possible).
\item Keep backups of every model.
\eit
\een

\chapter{Clean Code}\label{ch:clean_code}

Code is clean if it can be understood easily – by everyone on the team. Clean code can be read and enhanced by a
developer other than its original author. With understandability comes readability, changeability, extensibility and
maintainability.

\ben
\item \textbf{General Rules}
\bit
\item Follow standard conventions.
\item Keep it simple stupid. Simpler is always better. Reduce complexity as much as possible.
\item Boy scout rule. Leave the campground cleaner than you found it.
\item Always find root cause. Always look for the root cause of a problem.
\eit

\item \textbf{Design Rules}
\bit
\item Keep configurable data at high levels.
\item Prefer polymorphism to if/else or switch/case.
\item Separate multi-threading code.
\item Prevent over-configurability.
\item Use dependency injection.
\item Follow Law of Demeter. A class should know only its direct dependencies.
\eit

\item \textbf{Understandability Tips}
\bit
\item Be consistent. If you do something a certain way, do all similar things in the same way.
\item Use explanatory variables.
\item Encapsulate boundary conditions. Boundary conditions are hard to keep track of. Put the processing for them
in one place.
\item Prefer dedicated value objects to primitive type.
\item Avoid logical dependency. Don't write methods which works correctly depending on something else in the same class.
\item Avoid negative conditionals.
\eit

\item \textbf{Names Rules}
\bit
\item Choose descriptive and unambiguous names.
\item Make meaningful distinction.
\item Use pronounceable names.
\item Use searchable names.
\item Replace magic numbers with named constants.
\item Avoid encodings. Don't append prefixes or type information.
\eit

\item \textbf{Functions Rules}
\bit
\item Small.
\item Do one thing.
\item Use descriptive names.
\item Prefer fewer arguments.
\item Have no side effects.
\item Don't use flag arguments. Split method into several independent methods that can be called from the client
without the flag.
\eit

\item \textbf{Comments Rules}
\bit
\item Always try to explain yourself in code.
\item Don't be redundant.
\item Don't add obvious noise.
\item Don't use closing brace comments.
\item Don't comment out code. Just remove.
\item Use as explanation of intent.
\item Use as clarification of code.
\item Use as warning of consequences.
\eit

\item \textbf{Source Code Structure}
\bit
\item Separate concepts vertically.
\item Related code should appear vertically dense.
\item Declare variables close to their usage.
\item Dependent functions should be close.
\item Similar functions should be close.
\item Place functions in the downward direction.
\item Keep lines short.
\item Don't use horizontal alignment.
\item Use white space to associate related things and disassociate weakly related.
\item Don't break indentation.
\eit

\item \textbf{Objects \& Data Structures}
\bit
\item Hide internal structure.
\item Prefer data structures.
\item Avoid hybrids structures (half object and half data).
\item Should be small.
\item Do one thing.
\item Small number of instance variables.
\item Base class should know nothing about their derivatives.
\item Better to have many functions than to pass some code into a function to select a behavior.
\item Prefer non-static methods to static methods.
\eit

\item \textbf{Tests}
\bit
\item One assert per test.
\item Readable.
\item Fast.
\item Independent.
\item Repeatable.
\eit

\item \textbf{Code Smells}
\bit
\item Rigidity. The software is difficult to change. A small change causes a cascade of subsequent changes.
\item Fragility. The software breaks in many places due to a single change.
\item Immobility. You cannot reuse parts of the code in other projects because of involved risks and high effort.
\item Needless Complexity.
\item Needless Repetition.
\item Opacity. The code is hard to understand.
\eit
\een

\chapter{Software Engineering At Google}\label{ch:se_google}

\ben
\item \textbf{What Is Software Engineering?}
\bit
\item ``Software engineering'' differs from ``programming'' in dimensionality: programming is about producing code.
Software engineering extends that to include the maintenance of that code for its useful life span.
\item There is a factor of at least 100,000 times between the life spans of short-lived code and long-lived code. It
is silly to assume that the same best practices apply universally on both ends of that spectrum.
\item Software is sustainable when, for the expected life span of the code, we are capable of responding to changes
in dependencies, technology, or product requirements. We may choose to not change things, but we need to be capable.
\item Hyrum's Law: with a sufficient number of users of an API, it does not matter what you promise in the contract:
all observable behaviors of your system will be depended on by somebody.
\item Every task your organization has to do repeatedly should be scalable (linear or better) in terms of human input.
Policies are a wonderful tool for making process scalable.
\item Process inefficiencies and other software-development tasks tend to scale up slowly. Be careful about
boiled-frog problems.
\item Expertise pays off particularly well when combined with economies of scale.
\item ``Because I said so'' is a terrible reason to do things.
\item Being data driven is a good start, but in reality, most decisions are based on a mix of data, assumption,
precedent, and argument. It's best when objective data makes up the majority of those inputs, but it can rarely be
all of them.
\item Being data driven over time implies the need to change directions when the data changes (or when assumptions
are dispelled). Mistakes or revised plans are inevitable.
\eit


\item \textbf{How To Work Well On Teams}
\bit
\item Be aware of the trade-offs of working in isolation.
\item Acknowledge the amount of time that you and your team spend communicating and in interpersonal conflict. A
small investment in understanding personalities and working styles of yourself and others can go a long way toward
improving productivity.
\item If you want to work effectively with a team or a large organization, be aware of your preferred working style
and that of others.
\eit

\item \textbf{Knowledge Sharing}
\bit
\item Psychological safety is the foundation for fostering a knowledge-sharing environment.
\item Start small: ask questions and write things down.
\item Make it easy for people to get the help they need from both human experts and documented references.
\item At a systemic level, encourage and reward those who take time to teach and broaden their expertise beyond just
themselves, their team, or their organization.
\item There is no silver bullet: empowering a knowledge-sharing culture requires a combination of multiple
strategies, and the exact mix that works best for your organization will likely change over time.
\eit

\item \textbf{Engineering for Equity}
\bit
\item Bias is the default.
\item Diversity is necessary to design properly for a comprehensive user base.
\item Inclusivity is critical not just to improving the hiring pipeline for underrepresented groups, but to providing
a truly supportive work environment for all people.
\item Product velocity must be evaluated against providing a product that is truly useful to all users. It's better
to slow down than to release a product that might cause harm to some users.
\eit

\item \textbf{How To Lead A Team}
\bit
\item Don't ``manage'' in the traditional sense; focus on leadership, influence, and serving your team.
\item Delegate where possible; don't DIY (Do It Yourself).
\item Pay particular attention to the focus, direction, and velocity of your team.
\eit

\item \textbf{Leading At Scale}
\bit
\item Always Be Deciding: Ambiguous problems have no magic answer; they're all about finding the right trade-offs of
the moment, and iterating.
\item Always Be Leaving: Your job, as a leader, is to build an organization that automatically solves a class of
ambiguous problems—over time—without you needing to be present.
\item Always Be Scaling: Success generates more responsibility over time, and you must proactively manage the scaling
of this work in order to protect your scarce resources of personal time, attention, and energy.
\eit

\item \textbf{Measuring Engineering Productivity}
\bit
\item Before measuring productivity, ask whether the result is actionable, regardless of whether the result is
positive or negative. If you can't do anything with the result, it is likely not worth measuring.
\item Select meaningful metrics using the GSM framework. A good metric is a reasonable proxy to the signal you're
trying to measure, and it is traceable back to your original goals.
\item Select metrics that cover all parts of productivity (QUANTS). By doing this, you ensure that you aren't
improving one aspect of productivity (like developer velocity) at the cost of another (like code quality).
\item Qualitative metrics are metrics, too! Consider having a survey mechanism for tracking longitudinal metrics
about engineers' beliefs. Qualitative metrics should also align with the quantitative metrics; if they do not, it is
likely the quantitative metrics that are incorrect.
\item Aim to create recommendations that are built into the developer workflow and incentive structures. Even though
it is sometimes necessary to recommend additional training or documentation, change is more likely to occur if it is
built into the developer's daily habits.
\eit

\item \textbf{Style Guides And Rules}
\bit
\item Rules and guidance should aim to support resilience to time and scaling.
\item Know the data so that rules can be adjusted.
\item Not everything should be a rule.
\item Consistency is key.
\item Automate enforcement when possible.
\eit


\item \textbf{Code Review}
\bit
\item Code review has many benefits, including ensuring code correctness, comprehension, and consistency across a
codebase.
\item Always check your assumptions through someone else; optimize for the reader.
\item Provide the opportunity for critical feedback while remaining professional.
\item Code review is important for knowledge sharing throughout an organization.
\item Automation is critical for scaling the process.
\item The code review itself provides a historical record.
\eit

\item \textbf{Documentation}
\bit
\item Documentation is hugely important over time and scale.
\item Documentation changes should leverage the existing developer workflow.
\item Keep documents focused on one purpose.
\item Write for your audience, not yourself.
\eit

\item \textbf{Testing Overview}
\bit
\item Automated testing is foundational to enabling software to change.
\item For tests to scale, they must be automated.
\item A balanced test suite is necessary for maintaining healthy test coverage.
\item ``If you liked it, you should have put a test on it.''
\item Changing the testing culture in organizations takes time.
\eit


\item \textbf{Unit Testing}
\bit
\item Strive for unchanging tests.
\item Test via public APIs.
\item Test state, not interactions.
\item Make your tests complete and concise.
\item Test behaviors, not methods.
\item Structure tests to emphasize behaviors.
\item Name tests after the behavior being tested.
\item Don't put logic in tests.
\item Write clear failure messages.
\item Follow DAMP over DRY when sharing code for tests.
\eit

\item \textbf{Test Doubles}
\bit
\item A real implementation should be preferred over a test double.
\item A fake is often the ideal solution if a real implementation can't be used in a test.
\item Overuse of stubbing leads to tests that are unclear and brittle.
\item Interaction testing should be avoided when possible: it leads to tests that are brittle because it exposes
implementation details of the system under test.
\eit

\item \textbf{Larger Testing}
\bit
\item Larger tests cover things unit tests cannot.
\item Large tests are composed of a System Under Test, Data, Action, and Verification.
\item A good design includes a test strategy that identifies risks and larger tests that mitigate them.
\item Extra effort must be made with larger tests to keep them from creating friction in the developer workflow.
\eit


\item \textbf{Deprecation}
\bit
\item Software systems have continuing maintenance costs that should be weighed against the costs of removing them.
\item Removing things is often more difficult than building them to begin with because existing users are often using
the system beyond its original design.
\item Evolving a system in place is usually cheaper than replacing it with a new one, when turndown costs are included.
\item It is difficult to honestly evaluate the costs involved in deciding whether to deprecate: aside from the direct
maintenance costs involved in keeping the old system around, there are ecosystem costs involved in having multiple
similar systems to choose between and that might need to interoperate. The old system might implicitly be a drag on
feature development for the new. These ecosystem costs are diffuse and difficult to measure. Deprecation and removal
costs are often similarly diffuse.
\eit

\item \textbf{Version Control And Branch Management}
\bit
\item Use version control for any software development project larger than ``toy project with only one developer that
will never be updated.''
\item There's an inherent scaling problem when there are choices in ``which version of this should I depend upon?''
\item One-Version Rules are surprisingly important for organizational efficiency. Removing choices in where to commit
or what to depend upon can result in significant simplification.
\item In some languages, you might be able to spend some effort to dodge this with technical approaches like shading,
separate compilation, linker hiding, and so on. The work to get those approaches working is entirely lost labor—your
software engineers aren't producing anything, they're just working around technical debts.
\item Previous research (DORA/State of DevOps/Accelerate) has shown that trunk- based development is a predictive
factor in high-performing development organizations. Long-lived dev branches are not a good default plan.
\item Use whatever version control system makes sense for you. If your organization wants to prioritize separate
repositories for separate projects, it's still probably wise for interrepository dependencies to be unpinned/``at
head''/``trunk based.'' There are an increasing number of VCS and build system facilities that allow you to have both
small, fine-grained repositories as well as a consistent ``virtual'' head/trunk notion for the whole organization. \eit

\item \textbf{Code Search}
\bit
\item Helping your developers understand code can be a big boost to engineering productivity. At Google, the key tool
for this is Code Search.
\item Code Search has additional value as a basis for other tools and as a central, standard place that all
documentation and developer tools link to.
\item The huge size of the Google codebase made a custom tool—as opposed to, for example, grep or an IDE's
indexing—necessary.
\item As an interactive tool, Code Search must be fast, allowing a ``question and answer'' workflow. It is expected
to have low latency in every respect: search, browsing, and indexing.
\item It will be widely used only if it is trusted, and will be trusted only if it indexes all code, gives all
results, and gives the desired results first. However, earlier, less powerful, versions were both useful and used, as
long as their limits were understood.
\eit

\item \textbf{Build Systems And Build Philosophy}
\bit
\item A fully featured build system is necessary to keep developers productive as an organization scales.
\item Power and flexibility come at a cost. Restricting the build system appropriately makes it easier on developers.
\item Build systems organized around artifacts tend to scale better and be more reliable than build systems organized
around tasks.
\item When defining artifacts and dependencies, it's better to aim for fine-grained modules. Fine-grained modules are
better able to take advantage of parallelism and incremental builds.
\item External dependencies should be versioned explicitly under source control. Relying on ``latest'' versions is a
recipe for disaster and unreproducible builds.
\eit

\item \textbf{Critique: Google's Code Review Tool}
\bit
\item Trust and communication are core to the code review process. A tool can enhance the experience, but it can't
replace them.
\item Tight integration with other tools is key to great code review experience.
\item Small workflow optimizations, like the addition of an explicit ``attention set,'' can increase clarity and
reduce friction substantially.
\eit

\item \textbf{Static Analysis}
\bit
\item Focus on developer happiness. We have invested considerable effort in building feedback channels between
analysis users and analysis writers in our tools, and aggressively tune analyses to reduce the number of false
positives.
\item Make static analysis part of the core developer workflow. The main integration point for static analysis at
Google is through code review, where analysis tools provide fixes and involve reviewers. However, we also integrate
analyses at additional points (via compiler checks, gating code commits, in IDEs, and when browsing code).
\item Empower users to contribute. We can scale the work we do building and maintaining analysis tools and platforms
by leveraging the expertise of domain experts. Developers are continuously adding new analyses and checks that make
their lives easier and our codebase better.
\eit

\item \textbf{Dependency Management}
\bit
\item Prefer source control problems to dependency management problems: if you can get more code from your
organization to have better transparency and coordination, those are important simplifications.
\item Adding a dependency isn't free for a software engineering project, and the complexity in establishing an
``ongoing'' trust relationship is challenging. Importing dependencies into your organization needs to be done
carefully, with an understanding of the ongoing support costs.
\item A dependency is a contract: there is a give and take, and both providers and consumers have some rights and
responsibilities in that contract. Providers should be clear about what they are trying to promise over time.
\item SemVer is a lossy-compression shorthand estimate for ``How risky does a human think this change is?'' SemVer
with a SAT-solver in a package manager takes those estimates and escalates them to function as absolutes. This can
result in either overconstraint (dependency hell) or underconstraint (versions that should work together that don't).
\item By comparison, testing and CI provide actual evidence of whether a new set of versions work together.
\item Minimum-version update strategies in SemVer/package management are higher fidelity. This still relies on humans
being able to assess incremental version risk accurately, but distinctly improves the chance that the link between
API provider and consumer has been tested by an expert.
\item Unit testing, CI, and (cheap) compute resources have the potential to change our understanding and approach to
dependency management. That phase-change requires a fundamental change in how the industry considers the problem of
dependency management, and the responsibilities of providers and consumers both.
\item Providing a dependency isn't free: ``throw it over the wall and forget'' can cost you reputation and become a
challenge for compatibility. Supporting it with stability can limit your choices and pessimize internal usage.
Supporting without stability can cost goodwill or expose you to risk of important external groups depending on
something via Hyrum's Law and messing up your ``no stability'' plan.
\eit

\item \textbf{Large-Scale Changes}
\bit
\item An LSC process makes it possible to rethink the immutability of certain technical decisions.
\item Traditional models of refactoring break at large scales.
\item Making LSCs means making a habit of making LSCs.
\eit

\item \textbf{Continuous Integration}
\bit
\item A CI system decides what tests to use, and when.
\item CI systems become progressively more necessary as your codebase ages and grows in scale.
\item CI should optimize quicker, more reliable tests on presubmit and slower, less deterministic tests on post-submit.
\item Accessible, actionable feedback allows a CI system to become more efficient.
\eit


\item \textbf{Continuous Delivery}
\bit
\item Velocity is a team sport: The optimal workflow for a large team that develops code collaboratively requires
modularity of architecture and near-continuous integration.
\item Evaluate changes in isolation: Flag guard any features to be able to isolate problems early.
\item Make reality your benchmark: Use a staged rollout to address device diversity and the breadth of the userbase.
Release qualification in a synthetic environment that isn't similar to the production environment can lead to late
surprises.
\item Ship only what gets used: Monitor the cost and value of any feature in the wild to know whether it's still
relevant and delivering sufficient user value.
\item Shift left: Enable faster, more data-driven decision making earlier on all changes through CI and continuous
deployment.
\item Faster is safer: Ship early and often and in small batches to reduce the risk of each release and to minimize
time to market.
\eit

\item \textbf{Compute As A Service}
\bit
\item Scale requires a common infrastructure for running workloads in production.
\item A compute solution can provide a standardized, stable abstraction and environment for software.
\item Software needs to be adapted to a distributed, managed compute environment.
\item The compute solution for an organization should be chosen thoughtfully to provide appropriate levels of
abstraction.
\eit
\een

\chapter{Thinking Fast And Slow} \label{ch:thinking_fast_and_slow}

\section*{Introduction}

\bit
\item We use resemblance as a simplifying heuristic to make difficult judgment, causing predictable biases in
predictions.
\item Social scientists in the 1970s broadly accepted that people are generally rational, and emotions such as fear,
affection, and hatred explain departures from rationality.
\item People tend to assess the relative importance of issues by the ease with which they are retrieved from memory,
(eg words starting with a letter VS words having the letter in third row, news that are mentioned a lot by the media,
etc)
\item Accurate intuitions of experts are better explained by the effects of prolonged practice than by heuristics.
\item Valid intuitions develop when experts have learned to recognize familiar elements in a new situation and to act
in a manner that is appropriate to it.
\item When faced with a difficult question, we often answer an easier one instead, usually without noticing the
substitution.
\item When intuition fails, because neither an expert solution nor a heuristic answer comes to mind, we resort to
slower, deliberate, and effortful thinking.
\eit

\section*{Part 1: Two Systems}

\subsection*{The Characters Of The Story}
\bit
\item System 1 operates automatically and quickly, with no effort and no voluntary control.
\item System 2 allocates attention to effortful mental activities. It's associated with the subjective experience of
agency, choice, and concentration.
\item All operations of System 2 require attention and are disrupted when attention is drawn away.
\item System 2 has some ability to change the way System 1 works, by programming the normally automatic functions of
attention and memory.
\item System 1 generates impressions, intuitions, intentions, and feelings. Those endorsed by System 2 turn into
beliefs and voluntary actions.
\item Most of what System 2 thinks and does originates in your System 1, but System 2 takes over when things get
difficult, and it normally has the last word.
\item One of the tasks of System 2 is to overcome the impulses of System 1. System 2 is in charge of self-control.
\item System 2 is too slow and inefficient to substitute for System 1. The best we can do is to recognize when
mistakes are likely, and to try harder to avoid significant mistakes when the stakes are high.
\eit

\subsection*{Attention And Effort}
\bit
\item Pupils are sensitive indicators of mental effort. The more System 2 exerts mental effort, the more they dilate.
\item We decide what to do, but we have limited control over the effort of doing it. The task at hand decides this.
\item Orienting and responding quickly to the gravest threats or most promising situations improved the chance of
survival.
\item When System 2 is in action for a specific task is very easy to miss other events. Eg: The gorilla experiment,
when we are counting number of balls exchanged it is very easy to miss the gorilla.
\item In the economy of action, effort is a cost, and the acquisition of skill is driven by balancing benefits and
costs. Laziness is in our nature.
\item Law of least effort: if there are many ways to obtain something, we will eventually chose the most effortless.
\item System 2 is the only one that can follow rules, compare objects on several attributes, and make deliberate
choices between objects.
\item A crucial capability of System 2 is that it can program memory to obey an instruction that overrides habitual
responses.
\item Multitasking is effortful. Time pressure is another driver. Any task that requires keeping several ideas in
mind simultaneously has the same hurried character.
\eit

\subsection*{The Lazy Controller}
\bit
\item Mihaly's flow is a state of effortless concentration so deep that people lose a sense of time, of themselves,
and of their problems. This flow separates the two forms of effort: Concentration on the task and the deliberate
control of attention.
\item People who are cognitively busy are more likely to yield to temptation, make selfish choices, use sexist
language, and make superficial judgments in social situations. Controlling thoughts and behaviours is one of the
tasks that System 2 performs.
\item If you exert self-control for a task, then you are less willing or able to exert self-control for a following
task. This is called ego depletion.
\item When you are thinking hard or exerting self-control, your blood glucose level drops. The implication is that
you can undo ego depletion by ingesting glucose.
\item When people believe a conclusion is true, they are also very likely to believe arguments that appear to support
it, even when those arguments are unsound.
\item Intelligence is not only the ability to reason; it is also the ability to find relevant material in memory and
to deploy attention when needed.
\item System 2 is lazy. When an intuitive answer seems correct, we must deliberately put it in action to resist the
intuitive answer. However there are people that do make the extra step to think deeply (engaged) and people who don’t
(not engaged).
\item ``Engaged'' people are more alert, less willing to be satisfied with superficially attractive answers, more
skeptical about their intuitions., and more rational. People who are not ``engaged'' goes a lot with their emotions and
intuitions, are impulsive, impatient, and keen to receive immediate gratification.
\eit

\subsection*{The Associative Machine}
\bit
\item The responses by System 1 are associatively coherent, yielding a self-reinforcing pattern of cognitive,
emotional, and physical responses.
\item Cognition is embodied. You think with your body, not only with your brain.
\item Ideas are nodes in a vast network called associative memory, where causes link to effects, things to their
properties, and things to their categories.
\item Priming is a technique whereby exposure to one stimulus influences a response to a subsequent stimulus, without
conscious guidance or intention. For example, the word NURSE is recognised more quickly following the word DOCTOR
than following the word BREAD.
\item Priming is not restricted to concepts and words. Events that you are not even aware of prime your actions and
emotions.
\item The influence of an action by the idea is called the ideomotor phenomenon effect. For example, thinking of old
age makes you act old, and vice versa.
\item Money primes individualism: a reluctance to be involved with others, to depend on others, or to accept demands
from others.
\item The Lady Macbeth effect is a priming effect said to occur when response to a cleaning cue is increased after
having been induced by a feeling of shame
\item Reality consists largely of the story that your System 2 tells itself about what is going on. Priming phenomena
arise in System 1, and you have no conscious access to them. System 1 provides the impressions that often turn into
your beliefs, and is the source of impulses that often become the choices of our actions.
\eit

\subsection*{Cognitive Ease}
\bit
\item When things are going well we experience cognitive ease. When something is wrong we experience cognitive strain.
\item Cognitive strain is affected by both the current level of effort and the presence of unmet demands. This
mobilizes System 2.
\item A repeated experience, clear display, primed idea, and good mood all increase cognitive ease. This in turn
makes things feel familiar, true, good, and effortless.
\item It’s likely to assume familiarity with a word that we have recently encountered, even if we don’t remember when.
\item When strained, you are vigilant, suspicious, invest more effort, feel less comfortable, and make fewer errors.
But you are less intuitive and less creative.
\item Predictable illusions occur if judgment is based on an impression of cognitive ease or strain. For example,
frequent repetition makes people believe lies.
\item To craft a persuasive message, use high quality paper, bright colours, simple words, memorable verse with
rhymes, and quote the source with the simpler name.
\item Cognitive strain, whatever the source, mobilises System 2, which is more likely to reject the intuitive answer
suggested by System 1.
\item The mere exposure effect links the repetition of an arbitrary stimulus and the mild affection that people have
for it. It's stronger for stimuli that we don't consciously see.
\item Mood affects the operation of System 1. When we are uncomfortable and unhappy, we lose touch with our intuition.
\eit

\subsection*{Norms, Surprises, And Causes}

\bit
\item The main function of System 1 is to maintain and update a model of your personal world, which represents what
is normal in it.
\item Norm theory is when a surprising event happens, and subsequent similar surprising events will appear more
normal because they are interpreted in conjunction with it.
\item We have norms for a vast number of categories, which provide the background for the immediate detection of
anomalies.
\item System 1 is adept at finding a coherent causal story that links the fragments of knowledge at its disposal.
\item We are ready from birth to have impressions of causality, which do not depend on reasoning about patterns or
causation. They are products of System 1.
\item We are prone to apply causal thinking (find a reason) to situations that require statistical reasoning (admit
it taw just luck), but System 1 cannot do this, and System 2 requires necessary training.
\eit

\subsection*{A Machine For Jumping To Conclusions}

\bit
\item Jumping to conclusions is efficient if the jump saves time and effort, the conclusion is likely correct, and
the cost of an occasional mistake is acceptable.
\item System 1 bets on answers, where recent events and the current context have the most weight in determining and
interpretation. Otherwise, more distant memories govern.
\item In order to understand something we must first believe it.
\item System 1 is gullible and biased to believe, while System 2 is in charge of doubting and unbelieving, but System
2 is sometimes busy and often lazy.
\item Unlike scientists, which test hypotheses by trying to refute them, we seek data that are likely to be
compatible with the beliefs that we currently hold.
\item The tendency to like or dislike everything about a person, including things you have not observed, is known as
the halo effect.
\item The halo effect increases the weight of first impressions, sometimes to the point that subsequent information
is mostly wasted.
\item To derive the most useful information from multiple sources of evidence, you should always try to make these
sources independent of each other.
\item Sometimes individuals perform poorly while the group performs well. This is due to independent observations
which lead to uncorrelated errors that follow a normal distribution with zero mean, which means that on average the
error is zero.
\item The standard practice of open discussion gives too much weight to the opinions of those who speak early and
assertively, causing others to line up behind them.
\item System 1 excels at constructing the best possible story that incorporates ideas currently activated, but it
cannot allow for information that it does not have.
\item Jumping to conclusions facilitates the achievement of coherence and of the cognitive ease that causes us to
accept a statement as true. It explains overconfidence. Usually we neglect base rates and other phenomena.
\eit

\subsection*{How Judgments Happen}

\bit
\item System 1 continually assesses the problems that an organism must solve to survive. We equate good mood and
cognitive ease with safety and familiarity.
\item Faces with a strong chin and a slight confident-appearing smile exude confidence.
\item A judgment heuristic is falling back on a simpler assessment that is made quickly and automatically and is
available when System 2 must make its decision.
\item Because System 1 represents categories by a prototype or a set of typical exemplars, it deals well with
averages but poorly with sums.
\item System 1 allows matching intensity across diverse and unrelated dimensions. This mode of prediction by matching
is statistically wrong, although acceptable to both systems.
\item The control over intended computations is far from precise, and we often compute much more than we want or need.
This is called the mental shotgun.
\eit

\subsection*{Answering An Easier Question}
\bit
\item If we can't satisfactorily answer a hard target question, then System 1 invokes substitution by recalling and
answering an easier heuristic question.
\item After answering a heuristic question, System 1 uses intensity matching to translate this answer to an answer of
the target question.
\item The dominance of conclusions over arguments is most pronounced when emotions are involved.
\item An affect heuristic is when people let their likes and dislikes determine their beliefs about the world.
\item While self-criticism is one of the functions of System 2, it is more of an apologist for than a critic of the
emotions of System 1.
\eit

\section*{Part 2: Heuristics and Biases}

\subsection*{The Law Of Small Numbers}
\bit
\item System 1 is inept when faced with statistical facts, which change the probability of outcomes but do not cause
them to happen.
\item Extreme outcomes, both high and low, are most likely to be found in small than in large samples.
\item Even statistical experts pay insufficient attention to sample size, and have poor intuitions of sampling effects.
\item System 2 is capable of doubt, but sustaining doubt is harder work than sliding into certainty.
\item System 1 runs ahead of the facts and constructs a rich image based on scraps of evidence, causing us to
exaggerate the consistency and coherence of what we see.
\item The associative machine seeks causes. But instead of focusing on how the event came to be, the statistical view
relates it to what could have happened instead.
\item We do not expect to see regularity produced by a random process. When we detect what appears to be a rule, we
quickly reject the idea that the process is truly random.
\item We always translate patterns to causality and treat patterns as an opposed force to randomness.
\item We pay more attention to the content of messages than to information about their reliability. Consequently we
view the world in a simpler and more coherent way than the data justify.
\eit

\subsection*{Anchors}

\bit
\item An anchoring effect occurs when people consider a particular value for an unknown quantity before estimating
that quantity.
\item Adjusting your estimate away from the anchor is an effortful activity. Insufficient adjustment, where we accept
the anchor, is a sign of a weak or lazy System 2.
\item Anchoring is also a priming effect, which selectively evokes compatible evidence. This is the automatic
operation of System 1.
\item The anchoring index is 100\% for people who adopt the anchor as an estimate, and zero for people who are able
to ignore the anchor altogether.
\item Anchors that are obviously random can be just as effective as potentially informative anchors.
\item When negotiating, don't make an outrageous counteroffer to an outrageous proposal, but make a scene and make it
clear that you won't continue with that number on the table.
\item To resist anchoring effects, search your memory for arguments against the anchor. This negates the biased
recruitment of thoughts that produces these effects.
\item System 2 is susceptible to the biasing effect of some anchors that makes some information easier to retrieve.
It has control over or knowledge of the effect.
\eit

\subsection*{The Science Of Availability}

\bit
\item The availability heuristic replaces estimating the size of a category or frequency of an event with the ease
with which instances come to mind.
\item Salient events, dramatic events, and personal experiences versus experiences by others bias the ease with which
instances come to mind. This explains why everyone in a group may feel as though he or she does more than his or her
fair share.
\item By asking people to provide more instances of a given behaviour, you increase their struggle, and consequently
they conclude that they don't adopt that behaviour.
\item Judgment is influenced more by the ease of retrieval than the number of instances retrieved. Increasing the
number of requested instances therefore weakens judgment.
\item When you provide a spurious reason for the difficulty of retrieving a large number of instances, judgment is
again strengthened. The surprise is eliminated.
\item People who are personally involved in the judgment are more likely to consider the number of instances and less
likely to go by fluency.
\item Fluency of instances is a System 1 heuristic, which is replaced by a focus on content when System 2 engages.
\item The conclusion is that the ease with which instances come to mind is a System 1 heuristic, which is replaced by
a focus on content when System 2 is more engaged. Multiple lines of evidence converge on the conclusion that people
who let themselves be guided by System 1 are more strongly susceptible to availability biases than others who are in
a state of higher vigilance. The following are some conditions in which people ``go with the flow” and are affected
more strongly by ease of retrieval than by the content they retrieved: - When they are engaged in another effortful
task at the same time. - When they are in a good mood because they just thought of a happy episode in their life. -
If they score low on a depression scale. - If they are knowledgeable novices on the topic of the task, in contrast to
true experts. - When they score high on a scale of faith in intuition. - If they are (or are made to feel) powerful.
\eit

\subsection*{Availability, Emotion, And Risk}

\bit
\item Our expectations about the frequency of events are distorted by the prevalence and emotional intensity of the
messages to which we are exposed.
\item An availability cascade is the self-sustaining chain of events through which a minor event leads to public
panic and large-scale government action.
\item Policy is about what people want and what is best for them. The availability cascade is the mechanism through
which biases flow into policy.
\item When dealing with small risks, we either ignore them or give them far too much weight, with no middle ground.
\item Availability cascades may have a long-term benefit of calling attention to classes of risks and by increasing
the risk-reduction budget. \eit

\subsection*{Tom W's Speciality}

\bit
\item The percentage of a particular class in a population is called the base rate of that class.
\item How much an instance conforms to the stereotype of a particular class is called the representativeness of that
instance.
\item To determine how likely an instance belongs to a class, we ignore the base rate of the class and focus on the
representativeness of the instance.
\item Probability by representativeness is more accurate than chance guesses, but neglecting base rate information
that points in another direction is a statistical sin.
\item When endorsing representativeness, System 1 will automatically process the available information as if it were
true, unless you decide immediately to reject it.
\item Bayesian statistics govern how we should adjust the base rates given an account of representativeness.
\eit

\subsection*{Linda: Less Is More}

\bit
\item A conjunction fallacy is when we judge a conjunction of two events to be more probable than one of the events
in a direct comparison.
\item The most coherent stories are not the most probable, but they are plausible. And we confuse the notions of
coherence, plausibility, and probability.
\item Consequently, adding details to a scenario can make it more persuasive, but less likely to come true.
\item When performing single evaluation instead of joint evaluation, the less is more principle evaluates a
collection of items by its average and not its sum.
\item The sum-like nature of a variable is less obvious for probability than something more enumerable like money.
\item A question phrased as ``how many'' makes you think of individuals, while ``what percentage'' does not. This
decreases the incidence of the conjunction fallacy.
\eit

\subsection*{Causes Trump Statistics}

\bit
\item Statistical base rates are facts about a population to which a case belongs, but they are not relevant to the
individual case.
\item Causal base rates change your view of how the case came to be.
\item We neglect statistical base rates, and we easily combine causal base rates with other case-specific information.
\item Resistance to stereotyping is a laudable moral position, but neglecting valid stereotypes inevitably leads to
suboptimal judgments.
\item System 1 can deal with stores in which the elements are causally linked, but it is weak in statistical reasoning.
\item Individuals feel relieved of responsibility when they know that others have heard the same request for help.
\item When the outcome surprises us, we are unwilling to deduce the particular from the general, but are willing to
infer the general from the particular.
\item Consequently, we are more likely to learn something by finding surprises in our own behaviour than by hearing
surprising facts about people in general.
\eit

\subsection*{Regression To The Mean}

\bit\item Regression to the mean is when poor performance is followed by improvement, and good performance is followed by
deterioration.
\item We tend to be nice to other people when they please us and nasty when they do not, we are statistically
punished for being nice and rewarded for being nasty.
\item The discrepancies between two trials does not need a causal explanation. Often luck explains why one is a
significant outlier.
\item The correlation coefficient between two measures is a measure of the relative weight of the factors they share.
\item Whenever the correlation between two scores is imperfect, there will be regression to the mean.
\item Causal explanations will be evoked when we detect regression, but they will be wrong because regression to the
mean has an explanation but does not have a cause.
\eit

\subsection*{Taming Intuitive Predictions}

\bit
\item We are capable of rejecting information as irrelevant and false, but adjusting for smaller weaknesses in the
evidence is not something System 1 can do.
\item If we are asked for a prediction but substitute an evaluation of the evidence, we generate biased predictions
that completely ignore regression to the mean.
\item To create an unbiased prediction, start with a baseline estimate and an estimate derived from the evidence, and
choose the value between them that is proportional to your estimate of correlation.
\item Such unbiased predictions make errors, but they are smaller and do not favour either higher or lower outcomes.
\item Unbiased predictions permit predicting rare or extreme cases only when the information is very good, so you'll
never have the satisfaction of calling an extreme case.
\item Unbiased predictions are less preferred when error has varying types and severity, and extreme cases must be
called correctly even if it accumulates error elsewhere.
\item Your intuitions will deliver predictions that are too extreme and you will be inclined to put far too much
faith into them.
\item ``Extreme predictions and a willingness to predict rare events from weak evidence are both manifestations of
System 1. It is natural for the associative machinery to match the extremeness of predictions to the perceived
extremeness of evidence on which it is based—this is how substitution works. And it is natural for System 1 to
generate overconfident judgments, because confidence, as we have seen, is determined by the coherence of the best
story you can tell from the evidence at hand. Be warned: your intuitions will deliver predictions that are too
extreme and you will be inclined to put far too much faith in them.”
\eit

\section*{Part 3: Overconfidence}

\subsection*{The Illusion Of Understanding}

\bit
\item We are always ready to interpret behaviour as a manifestation of general propensities and personality traits,
or causes that you readily match to effects.
\item In a story, many of the important events that involve choices tempts us to exaggerate the role of skill and
underestimate the part of luck.
\item When you adopt a new view of the world, you immediately lose much of your ability to recall what you used to
believe before your mind changed. This causes us to underestimate the extent to which we were surprise by past
events, also known as hindsight bias.
\item Hindsight bias leads us to asses the quality of a decision not by whether the process was sound but by whether
its outcome was good or bad.
\item The sense-making System 1 makes us see the world as more tidy, simple, predictable, and coherent than it really
is. So we think that we can predict the future.
\item The comparison of firms that have been more or less successful is to a significant extent a comparison between
firms that have been more or less lucky.
\eit

\subsection*{The Illusion Of Validity}

\bit
\item Declarations of high confidence tell you that an individual has constructed a coherent story in his or her
mind, not necessarily that the story is true.
\item The persistence of individual differences in achievement is the measure by which we confirm the existence of
skill in someone.
\item Facts that challenge basic assumptions, and thereby threaten our livelihood and self-esteem, are simply not
absorbed. The mind does not digest them.
\item People can maintain an unshakeable faith, however absurd, when they are surrounded by a community of
like-minded believers.
\item A person who acquires more knowledge develops an enhanced illusion of skill and becomes unrealistically
overconfident. They have many excuses ready when proven wrong.
\item Errors in prediction are inevitable because the world is unpredictable. And high subjective confidence is not
an indicator of accuracy.
\eit

\subsection*{Intuitions VS Formulas}

\bit
\item Roughly 60\% of studies have shown significantly better accuracy for algorithms in comparisons of clinical and
statistical predictions.
\item One reason experts may be inferior is that they try to be clever, think outside the box, and consider complex
combinations of features. This actually reduces validity.
\item Another reason is that when asked to evaluate the same information twice, we frequently give different answers.
Unnoticed stimuli can substantially influence System 1.
\item Assigning equal weights to all predictors is often superior than varying weights found by multiple regression,
because it's not affected by accidents of sampling.
\item Consequently, we can develop useful algorithms without any prior statistical research. And back of the envelope
judgments are often good enough.
\item The aversion to algorithms making decisions is rooted in the strong preference that many people have for the
natural over the synthetic or artificial.
\item Intuition adds value, but only after a disciplined collection of objective information and disciplined scoring
of separate traits.
\item When creating your own formula for an interview procedure, pick at most six dimensions, and develop a 1 to 5
scale for each one.
\eit

\subsection*{Expert Intuition: When Can We Trust It?}

\bit
\item In the recognition-primed decision model, System 1 comes up with a plan. System 2 simulates it. If it works,
it's implemented. Otherwise it's tweaked or discarded.
\item Emotional learning may be quick, but expertise takes time to develop because it requires building a large
collection of mini-skills.
\item Confidence does not imply truth. The associative machine suppresses doubt and evokes ideas that are compatible
with the current dominant story.
\item Skilled intuitions come from environments where the environment is sufficiently regular to be predictable, and
we can learn these regularities through prolonged practice.
\item Human learning is normally efficient. If a strong predictive cue exists, human observers are likely to find it,
given a sufficient opportunity to do so.
\item Algorithms perform better in noisy environments because they detect weakly valid cues, and they use such cues
consistently.
\item The unrecognized limits of professional skill help explain why experts are often overconfident.
\item Judgments that answer the wrong question can also be made with high confidence.
\eit

\subsection*{The Outside View}

\bit\item Confidentially collecting the judgment of each person in a group makes better use of the knowledge of its members.
\item Our inside view judgment is overly optimistic, while the outside view judgment rightly adjusts a baseline
prediction.
\item We routinely discard statistical information, such as that offered by the outside view, when it's incompatible
with personal impressions of a case.
\item A planning fallacy is unrealistically close to a best-case scenario, and could be improved by consulting the
statistics of similar cases.
\item To counter the planning fallacy, reference class forecasting uses distributional information to create a
baseline prediction. This adopts an outside view.
\item People often, but not always, take on risky projects because they are overly optimistic about the odds they
face.
\eit

\subsection*{The Engine Of Capitalism}

\bit
\item The people who have the greatest influence on the lives of others are likely to be optimistic and
overconfident, and to take more risks than they realise.
\item The optimistic risk taking by entrepreneurs contributes to the economic dynamism of society, even if most risk
takers end up disappointed.
\item We rate ourselves below average on any task we find difficult, and so we are overly optimistic about our
standing on any activity we do moderately well.
\item Entrepreneurs imagine a future where their actions determine the outcome of the firm, and the outcome is not
affected by competitors. This is because they know so little about these competitors. This competition neglect begets
more competitors than the market can profitably sustain, and so the average outcome is a loss.
\item A wide confidence interval is a confession of ignorance, which is not socially acceptable for someone who is
paid to be knowledgeable. Society force us to be overconfident.
\item Optimism is highly valued, both socially and in the market, and so we reward the providers of dangerously
misleading information more than we reward truth tellers.
\item Optimism contributes to resilience by defending one's self image, where we take credit for successes but little
blame for failures.
\eit

\section*{Part 4: Choices}

\subsection*{Bernoulli's Errors}

\bit
\item Choices between simple gambles provide a simple model that shares important features with the more complex
decisions that researchers actually aim to understand.
\item Expected utility theory defined axioms of rationality. Prospect theory examines why we deviate this theory
under risk.
\item A risk-averse decision maker will choose a sure thing that is less than the expected value of a gamble, paying
a premium to avoid uncertainty.
\item Bernoulli proposed that the psychological gamble isn't the weighted average of the monetary values, but the
weighted average of the utility of these outcomes.
\item Prospect theory states that happiness is determined by the recent change in wealth, and not its absolute value.
It defines a reference point for happiness.
\eit

\subsection*{Prospect Theory}

\bit
\item Your personal wealth does not determine your attitudes to gains and losses. We like winning and dislike losing.
And we dislike losing more than we like winning.
\item In financial outcomes, the reference point can be the status quo, what you expect, or what you feel entitled to.
Better outcomes are gains. Worse outcomes are losses.
\item A principle of diminishing sensitivity applies to changes in wealth, both for gains and losses.
\item We typically only accept a bet if its loss aversion ratio, or its ratio of gains to losses, is a minimum
between 1.5 and 2.5.
\item In mixed gambles, where both a gain and a loss are possible, loss aversion causes extremely risk-averse choices.
\item In bad choices, where a sure loss is compared to a larger loss that is merely possible, diminishing sensitivity
causes risk seeking.
\item But prospect theory cannot deal with disappointment. It also cannot deal with regret, such as losing the gamble
and foregoing the sure option.
\eit

\subsection*{The Endowment Effect}

\bit
\item Indifference curves assume that your utility depends entirely on the present situation, and that the evaluation
of a possible job does not depend on your current job.
\item In labor negotiations and bargaining, the reference point and loss aversion is well understood, but the
omission in most scenarios is theory-induced blindness.
\item Loss aversions says the disadvantages of a change loom larger than its advantages. This introduces a bias that
favours the status quo.
\item Under the endowment effect, owning a good increases its value. Your minimum selling price greatly exceeds your
maximum buying price, violating rational economic behaviour.
\item The endowment effect applies to goods ``for use'', or consumed or enjoyed. It doesn't apply to goods ``for
exchange'', or traded for other goods.
\item Loss aversion is built into the automatic evaluations of System 1. Consider the baby who holds on fiercely to a
toy and shows agitation when it is taken away.
\item Selling goods activates regions of the brain that are associated with disgust as pain. Buying activates these
areas if the price is too high.
\item Veteran traders are unaffected by the endowment effect. They ask ``How much do I want to have this good,
compared with other things I could have instead''?
\item Being poor is living below one's reference point. Small amounts of money they receive are a reduced loss, not a
gain. All their choices, then, are between losses.
\eit

\subsection*{Bad Events}

\bit
\item Our brain prioritises threats above opportunities. It can process hostile images that we can't consciously see,
and can quickly pick hostile faces out from a crowd. This happens because evolutionary speaking, we developed a
mechanism for detection threats abut not one for exploiting opportunities.
\item Even symbolic threats evoke in attenuated form many reactions of the real thing, including fractional
tendencies to avoid or approach, recoil or lean forward.
\item Given a goal, not achieving it is a loss, while exceeding it is a gain. But the aversion to the failure of not
reaching a goal is stronger than the desire to exceed it.
\item Negotiations are difficult because losses by one side induce pain that exceeds the pleasure of the gains by the
other side.
\item An existing wage, price, or rent sets a reference point that is entitled. We think that exploiting market power
to impose losses on others is unacceptable.
\item A firm has its own entitlement, which is to retain its current profit. But we think it's not unfair for a firm
to reduce its workers' wages when its profitability is falling.
\item If a merchant lowers the price of a good, customers who bought at the higher price think of themselves as
having sustained a loss that is more than appropriate.
\item Punishing one stranger for behaving unfairly to another stranger activates the pleasure center of the brain.
\eit

\subsection*{The Fourfold Pattern}

\bit
\item When mapping decision weights to outcome probabilities, they are not equal in value, contrary to the
expectation principle.
\item We overweight unlikely events, which illustrates the possibility effect. We underweight highly likely events,
which illustrates the certainty effect.
\item We also have inadequate sensitivity to intermediate probabilities. The range of decision weights is much
smaller than the range of probabilities.
\item Paying a premium to eliminate a worry with certainty is compatible with the psychology of worry but not with
the rational model.
\item Between a sure loss and a gamble with a high probability of a larger loss, diminishing sensitivity makes the
sure loss more aversive, and the certainty effect reduces the aversiveness of the gamble.
\item This explains why people accept a high probability of making things worse for a small hope of avoiding a large
loss, which can turn manageable failures into disasters.
\item These same two factors enhance the attractiveness of the sure thing and reduce the attractiveness of the gamble
when the out come is positive.
\item Systematic deviations from expected value are costly in the long run. This rule applies to both risk aversion
and to risk seeking.
\eit

\subsection*{Rare Events}

\bit
\item Emotion and vividness influence fluency, availability, and judgments of probability, and thus account for our
excessive response to the few rare events we don't ignore.
\item People overestimate the probabilities of unlikely events, and overweight unlikely events in their decisions.
\item If the event we are asked to estimate is very unlikely, we instead focus on its alternative. We focus on the
odd, different, and unusual.
\item A rich and vivid representation of the outcome, whether or not it is emotional, reduces the role of probability
in the evaluation of an uncertain prospect.
\item The denominator effect is when your attention is drawn to one outcome that ``stands out,'' and you do not assess
the alternative outcome with the same care.
\item We weight low-probability events more when stated in terms of relative frequencies (how many) than when stated
in more abstract terms of ``chances,'' ``risk,'' and ``probability'' (how likely).
\item In choice from experience, instead of choices from description, we are exposed to variable outcomes from the
same source, and so we do not overweight rare events.
\eit

\subsection*{Risk Policies}

\bit
\item Every simple choice formulated as gains and losses can be deconstructed in innumerable ways into a combination
of choices, yielding preferences that are likely to be inconsistent.
\item If paying a premium for sure gains and to avoid a sure loss come out of the same pocket, the discrepant
attitudes are unlikely to be optimal.
\item We prefer narrow framing, or considering a sequence of simple decisions, even when we must entertain broad
framing, or considering the decisions jointly.
\item Narrow framing of gambles leads to loss aversion. Broad framing treats each gamble as one of many, blunting the
emotional reaction to loss and increasing the tolerance of risk.
\item Closely following daily fluctuations is a losing proposition, because the pain of small losses exceeds the
pleasure of equally frequent small gains.
\item A risk policy eliminates the pain of occasional loss by the thought that the policy that left you exposed to it
will be advantageous to you over the long run.
\item While an outside view protects you from the exaggerated optimism of the planning fallacy, a risk policy
protects you from the exaggerated caution induced by loss aversion.
\eit

\subsection*{Keeping Score}

\bit
\item The emotion that people attach to the state of our mental accounts are not acknowledged in standard economic
theory.
\item The disposition effect is the bias in finance to sell winners rather than losers, and is an instance of narrow
framing.
\item The sunk-cost fallacy invests additional resources in a losing account when better investments are available.
It prefers an unfavorable gamble to a sure loss.
\item Members of a board will replace a CEO with one who does not carry the same mental accounts and is therefore
better able to ignore sunk costs of past involvements.
\item A poignant story evokes more regret if it involves unusual events, because such events attract attention and
are easier to undo in our imagination.
\item People expect to have stronger emotional reactions to an outcome that is produced by action than to the same
reaction when it is produced by inaction.
\item When you deviate from the default, you can easily imagine the norm. If the default is associated with bad
consequences, the discrepancy can be a source of painful emotions.
\item You will be more loss averse in situations that are more important than money, and more reluctant to sell
important endowments when it might lead to an awful outcome.
\item To inoculate yourself against regret, remind yourself of its possibility and that things can go badly, and
preclude any hindsight that might cause it.
\eit

\subsection*{Reversals}

\bit
\item We normally experience situations in which contrasting alternatives are absent, and so moral intuitions that
come to your mind in different scenarios are inconsistent.
\item Preference reversal can occur if joint evaluation focuses attention on a situational aspect that is less
salient than in single evaluation.
\item The emotional reactions of System 1 likely determine single evaluation, while the comparison and careful
evaluation required by joint evaluation calls for System 2.
\item Judgments and preferences are coherent within categories but potentially incoherent when the objects are
evaluated belonging to different categories.
\item The evaluability hypothesis states that some attributes are only given weight in a joint evaluation because
they are not evaluable on their own.
\item Be wary of joint evaluation when someone who controls what you see has a vested interest in what you choose.
\eit

\subsection*{Frames And Reality}

\bit
\item Losses evoke stronger negative feelings than costs do, and so the cost of a lottery ticket that did not win is
more acceptable than losing a gamble.
\item ``Rational'' subjects, which are least susceptible to framing effects, showed enhanced activity in the frontal
area of the brain that is implicated in combining emotion and reasoning.
\item Reframing is effortful and System 2 is normally lazy. Most of us passively accept decisions as they are framed,
and so our preferences are frame-bound and not reality-bound.
\item Broader frames and inclusive accounts generally lead to more rational decisions.
\item Opt-in versus opt-out is another framing effect. We check a box if we've already decided what we wish to do.
But if unprepared for the question, laziness prefers the default.
\eit

\section*{Part 5: Two Selves}

\subsection*{Two Selves}

\bit
\item The decision maker who pays different amounts to achieve the same gain or be spared the same loss is making a
mistake.
\item The peak-end rule states that a global retrospective rating is dominated by the highest score (the peak) and
the final score (the end).
\item The duration neglect states that the length of the evaluated activity has no effect whatsoever on its evaluation.
\item The experiencing self, which evaluates each moment, does not have a voice. The remembering self, which keeps
score and governs what we learn, makes our decisions.
\item What we learn from the past is to maximize the qualities of our future memories, not necessarily our future
experience.
\item This is a feature of System 1, which represents sets by averages, norms, and prototypes. Not by sums.
\item In some cases rats who can stimulate their brain by pressing a lever will die of starvation without taking a
break to feed themselves.
\item A memory that neglects duration will not serve our preference for long pleasures and short pains.
\eit

\subsection*{Life As A Story}

\bit
\item Duration neglect is normal in a story. A story is about significant events and memorable events like its
ending, not about time passing.
\item The peak-end rule and duration effect also influence our evaluation of others' lives.
\item Vacation pictures may be important to the remembering self. The photographer does not view the scene as a
moment to be savored but as a future memory to be designed.
\item We use the word ``memorable'' to describe vacation highlights, which explicitly reveals the goal of the
experience.
\eit

\subsection*{Experienced Well-Being}

\bit
\item A small fraction of the population endures most of the suffering, whether because of illness, an unhappy
temperament, or misfortunes or tragedies.
\item Our emotional state is largely determined by what we attend to, and we are normally focused on our current
activity and immediate environment.
\item Less commuting, greater availability of child care, and improved socializing opportunities can reduce the
``unhappiness index'' of a society.
\item From Gallup data, some life aspects, such as education, is associated with a higher evaluation of one's life,
but not with greater experienced well-being.
\item Religion has favorable impact on both positive affect and stress reduction than on life evaluation, but
provides no reduction of feelings of depression or worry.
\item Being poor makes one miserable. A salary beyond \$75,000 may enhance one's life satisfaction, but it does not
improve experienced well-being.
\item Higher income allows the purchase of many pleasures, but it also reduces the ability to enjoy the smaller
things in life. Therefore the emotional experience is unchanged.
\eit

\subsection*{Thinking About Life}

\bit
\item Affective forecasting is the forecast of one's personal state in the future. An error happens when you think
statistics don't apply to you.
\item The score you assign to your life is determined by a small sample of highly available ideas, and not a careful
weighing of the domains of your life.
\item Both experienced temperament and life satisfaction are largely determined by the genetics of temperament.
\item Setting goals that are especially difficult to attain leads to a dissatisfied adulthood. We must consider what
people want in the concept of well-being.
\item The focusing illusion states that nothing in life is as important as you think it is when you are thinking
about it.
\item This illusion can cause people to be wrong about their present state of well-being as well as the happiness of
others, and about their own happiness in the future.
\item The term miswsanting describes bad choices that arise from the errors of affective forecasting. The focusing
illusion is a rich source of it.
\item The focusing illusion favours goods or experiences that are initially exciting but may lose their appeal. It
appreciates less experiences that retain their attention value.
\eit